\documentclass[instructions.tex]{subfiles}
\begin{document}
 
\chapter{Combien l’injustice est commune.}
 
Les riches et les puissans du siècle disent tous qu’ils ne font tort à personne. Le peuple et les pauvres allèguent des excuses. Ceux-là disent que, parmi le peuple, la plupart sont des voleurs; et ceux-ci prétendent que les plus grands voleurs sont parmi les puissans et les riches : c’est à chacun de sonder sa conscience. Comme la vérité ne fait acception de personne, nous allons examiner l’un et l’autre.
 
\primo Combien de gens vont la tête levée et se croient innocens, parce qu’on ne sait pas leur fourberie ou qu’on n’ose les reprendre ! Tel aujourd’hui passe pour un honnête homme, qui devant Dieu est un grand larron. Pour couvrir sa mauvaise foi, il crie lui-même le premier à l’injustice et accuse les autres de friponnerie. Dès qu’un homme puissant aime l’intérêt, il se croit tout permis; s’il ne prend pas de vive force le bien et les droits d’autrui, il sait par la chicane et par la surprise en venir à bout. On ne fait pas toujours mourir le pauvre Naboth, mais on n’en prend pas moins sa vigne.
 
On n’ose souvent ni se plaindre ni demander justice. Un seigneur doit être le protecteur et le père de ses sujets; mais ce qu’il exige d’eux est-il toujours légitime ? Se sert-il de son crédit et de ses biens pour les soutenir et les aider ? Les grands se servent-ils de leur autorité pour secourir le pauvre peuple ? Riches, grands du monde, puissans de la terre, que vous serez sévèrement jugés sur l’acquisition de vos biens, sur l’exaction de vos droits, sur l’usage que vous faites de votre autorité\footnote{Judicium durissimum his qui præsumt fiet. Sap. 6.} !
 
La justice est ouverte à tous; mais comment en sortent ceux qui sont sans appui ? Sa balance ne penche-t-elle jamais du côté de l’argent ou de la faveur ? La chicane n’a-t-elle point d’accès au barreau ? L’ignorance et la prévention ne montent-elles jamais sur les tribunaux ? Jugemens de Dieu, que vous dévoilerez un jour de mystères !
 
Les puissans qui oppriment les autres, les riches qui profitent de tout pour amasser, sont des fléaux de Dieu. Ils sont comme des verges dont le Tout-Puissant se sert pour punir les peuples et éprouver ses élus; mais ces verges seront jetées au feu, après avoir servi aux vengeances divines.
 
\secundo Le peuple et les pauvres excusent leurs larcins et leurs friponneries sur la nécessité. Ne faudrait-il pas mieux les imputer à leur fainéantise et à leurs débauches ? Sans prévoyance, sans économie, ils dépensent quelquefois dans une semaine ce qui les nourrirait pendant un mois, Paresseux, allez à la fourmi; apprenez d’elle à travailler. Qui sait employer le temps, gagne toujours de quoi vivre. Quiconque ne veut pas travailler, dit saint Paul, ne mérite pas de manger.
 
Combien de gens, sous prétexte qu’ils sont pauvres, ne vivent que de rapines et de larcins, ne nourrissent leur famille et leur bétail que des dépouilles d’autrui ! Que d’ouvriers, de domestiques, d’artisans qui ne travaillent que servilement, cachent artificieusement leurs friponneries, et dont les ouvrages trompeurs font connaître la mauvaise foi !
 
On voit aujourd’hui ce que le Prophète a prédit, que le mensonge et le larcin ont inondé la terre. Tel rougirait de voler sur les grands chemins, qui ne rougit point de le faire en secret, de tirer intérêt du pur prêt, de faire le monopole, de frauder dans un partage, de mentir dans une vente, de faire revivre ou d’acheter des actions injustes, éteintes ou prescrites. Que de bornes transposées ! que de forêts impunément dégradées ! que de revenus publics et de revenus de fabriques dissipés ! que de comptes refusés ! que de parjures en justice ! que de procès entortillés dans la chicane ! que de transactions forcées ! que d’impôts injustement répartis ! que de titres retenus et cachés ! que de vols recélés !
 
Ne voit-on pas des pères frauder la légitime de leurs enfans par des contrats ou des testamens sans règle; des mères et des enfans voler leur famille; des maris refuser à leurs femmes une assurance de la dot, et les forcer à des donations ou à des dettes indiscretes; des maîtres cruels refuser le salaire et une suffisante nourriture à des ouvriers et à des domestiques qu’ils accablent de travaux; des sujets éluder avec malice les droits légitimes\footnote{Mendacium.... et furtum inundaverunt.... propter quod lugebit terra. Os. 4.} ? O que l’attachement aux biens, l’injustice et la mauvaise foi ont fait périr d’âmes !
 
\section{Résolutions}
 
\primo Je craindrai souverainement l’injustice. 

\secundo Je remplirai tous les devoirs de mon état avec toute l’exactitude dont je suis capable.

\tertio Et dans toutes mes affaires avec le prochain, je mettrai la droiture, la franchise et la bonne foi qu’exige l’équité la plus exacte. 

Mon Dieu, enracinez profondément dans mon cœur l’amour de la justice, et inspirez-moi une vive horreur pour tout ce que condamne cette vertu si nécessaire.
 
\end{document}
